% -----------------------------------------------------------------------------
% Final discussion of the MKP, CEC2022, and HRES2--H2 results.
% This is an input fragment, not a standalone LaTeX document.
% Required packages: amsmath, booktabs, graphicx, and array.
% -----------------------------------------------------------------------------

\section{Discussion}
\label{sec:discussion-final}

The experiments provide consistent evidence that the proposed
DTW/DDTW-driven rotational framework is a robust way to exploit a portfolio of
metaheuristics across discrete, continuous, and engineering optimization
problems (Tables~\ref{tab:final-mkp-wtl}, \ref{tab:final-cec-wtl}, and
\ref{tab:final-hres-tests}). Its main benefit is not universal dominance on every problem, but a
reduction in dependence on the choice of a single standalone solver. This
distinction is important. In MKP, the framework accumulated between 40 and 41
significant wins across 63 baseline comparisons, depending on the
configuration, but it also incurred between 3 and 8 significant losses. In
CEC2022, DTW and DDTW achieved the best aggregate ranks among the six evaluated
methods, although specialized baselines remained preferable on some functions.
In HRES2--H2, both variants produced lower mean LCOE than all ten standalone
methods and all distributional comparisons remained significant after Holm
correction. Taken together, these findings support the framework as an
effective portfolio controller, while ruling out the stronger claim that it
must outperform every constituent algorithm on every landscape.

\begin{table}[htbp]
\centering
\caption{Cross-domain interpretation of the principal empirical findings.}
\label{tab:discussion-synthesis}
\small
\resizebox{\textwidth}{!}{%
\begin{tabular}{p{0.12\textwidth}p{0.36\textwidth}p{0.34\textwidth}p{0.30\textwidth}}
\toprule
Domain & Main evidence & Interpretation & Scope boundary \\
\midrule
MKP & DTW--3K obtained mean rank 1.22; the 3K budget reduced the average gap by 19.3\% for DTW and 21.5\% for DDTW. & Longer search materially improved solution quality, and level-based DTW was the strongest tested configuration. & Nine selected instances and no fixed-schedule or no-monitor control. \\
CEC2022 & DTW and DDTW obtained mean ranks 2.323 and 2.358, with W/T/L totals of 38/16/6 and 37/19/4. & Both variants were competitive over heterogeneous continuous landscapes; neither alignment mode was uniformly preferable. & Dimension $D=10$, 12 functions, and five standalone baselines. \\
HRES2--H2 & Both variants reached mean LCOE 0.267160 CNY/kWh and W/T/L $=10/0/0$; direct DTW--DDTW $p=0.593$. & The portfolio reliably reached a narrow high-quality basin, but no DTW--DDTW difference was detected. & One fixed synthetic year and independent hybrid--baseline samples. \\
\bottomrule
\end{tabular}%
}
\end{table}

\subsection{Portfolio adaptation across domains}
\label{sec:discussion-cross-domain}

The three domains differ in representation, objective scale, constraint
structure, and landscape geometry. MKP is a discrete, constrained maximization
problem with plateaus and combinatorial neighborhoods; CEC2022 contains shifted,
rotated, hybrid, and composition functions in a continuous space; and HRES2--H2
is a low-dimensional but tightly constrained techno-economic sizing problem in
which every objective evaluation includes an 8,760-hour dispatch simulation.
The fact that the framework remained competitive in all three settings is
therefore more informative than a gain observed on a single benchmark family.
The use of trajectory shape and run-relative percentile thresholds appears to
make the switching rule portable across objective scales. This interpretation
is consistent with the observed results, although the present design does not
isolate the contribution of percentile normalization from the other framework
components.

The W/T/L profiles also clarify where the portfolio is most useful. In MKP, the
hybrid defeated PSO on all nine instances in every configuration and generally
performed strongly against EHO, ACO, GWO, and WOA. GA was the most difficult
competitor, causing three significant losses in each DTW campaign and three or
four in the DDTW campaigns. Thus, rotation does not erase the value of a
well-matched specialist: on some discrete instances, GA remains a better choice.
In CEC2022, the most consistent hybrid advantage was against WOA, whereas ACO
was the strongest baseline in terms of hybrid losses. These patterns suggest
that the framework's value is greatest when its constituent solvers provide
complementary search behavior, rather than when one solver is already
consistently well aligned with the problem.

\subsection{MKP: budget effect and DTW--3K}
\label{sec:discussion-mkp}

The MKP results in Tables~\ref{tab:mkp_dtw_1k}--\ref{tab:mkp_ddtw_3k}
provide the clearest discrimination among the four framework configurations.
The global Friedman test was significant
($\chi_F^2=24.0667$, $p=2.419\times10^{-5}$), and all six paired configuration
contrasts remained significant after Holm correction. DTW--3K obtained the
best mean rank (1.22) and the lowest average normalized gap (1.805\%), followed
by DDTW--3K (rank 1.78; gap 1.876\%), DTW--1K (rank 3.11; gap 2.237\%), and
DDTW--1K (rank 3.89; gap 2.391\%); the corrected pairwise results are reported
in Table~\ref{tab:final-mkp-config-tests}. DTW--3K also achieved the lowest mean gap in
seven of the nine selected instances, while DDTW--3K was best in
\texttt{mknapcb2} and \texttt{mknapcb9}. Accordingly, DTW--3K is the best
supported configuration for this MKP test set.

The budget effect is both statistically and practically relevant. Increasing
the search from 1,000 to 3,000 iterations reduced the average gap by 19.3\% for
DTW and 21.5\% for DDTW. This result indicates that the hybrid had not exhausted
its useful search dynamics at 1,000 iterations. The corresponding mean number
of switches increased from 7.85 to 26.70 for DTW and from 16.43 to 26.96 for
DDTW. However, switch frequency and solution quality should not be conflated:
the logs show that more switching occurred under the longer budget, but they do
not demonstrate that every individual rotation generated an improvement.

At 1K, DDTW switched approximately twice as often as DTW while obtaining a
higher mean gap. One plausible explanation is that derivative-based alignment
is more sensitive to short changes around discrete plateaus and therefore
triggers more rotations without necessarily improving the incumbent solution.
By contrast, the absolute level information retained by DTW may be useful when
fitness trajectories contain long plateaus separated by discrete jumps. This
is a mechanism-level hypothesis, not a causal conclusion: it should be tested
with event-level analyses linking each switch to subsequent improvement, and
with controlled experiments that vary the monitoring window and patience
parameters.

The improved 3K performance also entails a computational trade-off. The
observed mean campaign times increased from 306.94 to 756.11 seconds for DTW
and from 256.43 to 682.21 seconds for DDTW. These measurements combine the
larger evaluation budget, solver cost, machine conditions, and monitoring
overhead; consequently, they quantify the empirical quality--time trade-off of
the campaigns but do not isolate the cost of DTW or DDTW itself. A fair overhead
study would compare identical solver trajectories with and without the monitor
under controlled hardware and evaluation counts.

\subsection{CEC2022: competitive but not universal}
\label{sec:discussion-cec2022}

The CEC2022 experiments in Table~\ref{tab:cec2022_results} show that the
framework generalizes beyond the discrete setting. DTW and DDTW achieved the
best aggregate mean ranks among the six evaluated algorithms, 2.323 and 2.358,
respectively (Table~\ref{tab:final-cec-ranks}). Their Holm-adjusted W/T/L
totals were 38/16/6 for DTW and 37/19/4 for DDTW across 60
hybrid--baseline comparisons. All 12 function-wise Friedman tests were
significant in each campaign, confirming that algorithm choice affected the
observed objective values. The framework therefore offers broad competitive
coverage, but the losses and ties show that no single hybrid variant dominated
the entire suite.

This landscape dependence is also visible in the descriptive results. DTW
reported the best hybrid mean on F6, F7, F8, and F11, whereas DDTW did so on
F10 and F12; on other functions, a standalone method could retain the best
mean. At the best-run level, both variants reached the official optimum on F1,
F5, and F11. Best-run results demonstrate attainable performance but are more
sensitive to stochastic selection than distributions over 31 runs, so the
rank and paired-test results should carry greater inferential weight. The large
dispersion on F6 (standard deviations of approximately 1,821.985 for DTW and
1,942.164 for DDTW) further illustrates why isolated best values are
insufficient to characterize robustness.

The direct DTW--DDTW comparison found no Holm-significant difference on any of
the 12 functions. DTW had the lower unrounded mean on seven functions and a
slightly lower average normalized mean gap (10.281\% versus 10.501\%), but the
evidence does not support declaring one alignment modality generally superior
for CEC2022. Equally, a non-significant result is not evidence of equivalence.
An equivalence or non-inferiority design with a prespecified practically
relevant margin would be required to claim that the two variants are
interchangeable.

Interestingly, switching activity in CEC2022 was much lower than in the MKP
3K and HRES2--H2 campaigns: DTW averaged 6.20 switches and DDTW 6.01, with a
range of four to ten. Because the stopping budgets and search spaces differ,
these raw counts are not directly comparable as rates. Nevertheless, their
variation across functions and domains is consistent with a data-dependent
trigger rather than a fixed rotation timetable. A future analysis should
normalize events by eligible monitoring iterations and relate them to
post-switch improvement, diversity recovery, and solver identity.

\subsection{HRES2--H2: statistical evidence and engineering relevance}
\label{sec:discussion-hres2}

In the engineering case, both variants converged to almost the same economic
region (Table~\ref{tab:hres2_metrics}). Their archived mean LCOE was 0.267160 CNY/kWh, with best values of
0.267159 CNY/kWh. The associated mean LCOH values were 17.631382 CNY/kg for DTW
and 17.631312 CNY/kg for DDTW. The best designs also shared the same reported
capacity pattern: approximately 174.47 MW of wind, 25.53 MW of photovoltaic
generation, 70 MW of electrolyzers, and a 50 MW/4 h battery, with annual
hydrogen production of approximately $8.86\times10^6$ kg. The renewable-share
constraint was active at, or very near, its 20\% upper bound. This boundary
solution indicates that the optimizer exploited the allowable grid support to
minimize cost under the stated formulation.

Against the ten standalone methods, all baseline mean LCOE values were higher,
and all Mann--Whitney comparisons in Table~\ref{tab:final-hres-tests} remained significant after Holm correction,
yielding W/T/L $=10/0/0$ for each framework variant. The magnitude of the
improvement was nevertheless heterogeneous. The relative reduction in mean
LCOE ranged from approximately 0.009\% against ILS to 6.73\% against SA for
DTW, and from approximately 0.009\% against ILS to 7.90\% against SA for
DDTW. Therefore, statistical significance should not be interpreted as uniform
engineering importance. Against ILS and TS, the absolute advantage is very
small and becomes statistically detectable because the distributions are
highly concentrated; against weaker baselines such as SA, the economic
difference is materially larger.

The direct same-seed comparison between DTW and DDTW produced a mean LCOE
difference of only $7.419\times10^{-7}$ CNY/kWh and a two-sided Wilcoxon
$p=0.593$. The data therefore do not detect a performance difference between
the alignment variants in this case. They do not establish formal equivalence,
particularly because objective values were archived at six decimal places and
the between-run dispersion was extremely small. The narrow distributions may
reflect stable behavior of the framework, but they may also arise from the
low-dimensional constrained basin, the deterministic dispatch rules, and the
use of one fixed weather and load profile.

The HRES2--H2 interpretation must also remain aligned with the implemented
system boundary. The case models grid-connected wind, photovoltaic generation,
an electrolyzer, and a battery for annual hydrogen production. It does not
include a hydrogen storage tank, fuel-cell reconversion, or a minimum hydrogen
demand constraint. Consequently, the results validate the optimizer for the
specified wind--PV--electrolyzer--battery sizing and dispatch problem, not for
every configuration commonly described as a complete hydrogen-storage
microgrid. Robust engineering conclusions will require multiple weather and
price years, demand scenarios, technology-cost trajectories, and uncertainty
or reliability constraints.

\subsection{What can be inferred from the switching dynamics}
\label{sec:discussion-switching}

The switch counts differed substantially by domain and budget: 7.85 and 16.43
for DTW--1K and DDTW--1K in MKP; 26.70 and 26.96 under the 3K budget; 6.20 and
6.01 in CEC2022; and 21.00 and 21.68 in HRES2--H2. This heterogeneity is
compatible with the intended behavior of an adaptive controller whose trigger
depends on the recent trajectory rather than on a fixed number of iterations.
The similarity between DTW and DDTW in CEC2022 and HRES2--H2 also suggests that
both representations identified comparable stagnation structure in those
campaigns, whereas their larger separation in MKP--1K points to a stronger
effect of trajectory representation on discrete plateaus.

The synchronized HRES2--H2 convergence and diagnostic profiles in
Figures~\ref{fig:hres2_dtw_conv}--\ref{fig:hres2_ddtw_delta} visually connect
fitness plateaus with the monitor signals and solver handoffs. These plots are
useful for explaining the operating principle, but they correspond to selected
best-run traces rather than an average trajectory over 31 runs. They should
therefore be interpreted as mechanism illustrations, not as independent
evidence that the same temporal sequence occurred in every run.

These counts remain diagnostic rather than causal evidence. They do not state
which solver was entered or left, whether diversity increased, how long the
benefit persisted, or whether an improvement would have occurred without the
switch. A stronger causal analysis would compare the observed policy with
fixed-period rotation, random rotation, a no-switch portfolio, and the same
solvers run independently under equal evaluation budgets. Event studies could
then estimate the change in improvement rate immediately before and after a
trigger and identify which transitions contribute most to performance.

\subsection{Role of shared elitist memory}
\label{sec:discussion-memory}

The framework transfers the incumbent elite solution across solver rotations.
This design is important because it protects the best solution found so far
while allowing the active search dynamics to change. It offers a plausible
explanation for why rotation can diversify the search without sacrificing
accumulated progress, especially in MKP and HRES2--H2 where returning to a
previously identified high-quality region is valuable. At the same time,
aggressive elitist transfer can reduce effective diversity if each solver is
reinitialized too close to the same incumbent. The current results demonstrate
the performance of the complete framework but do not separate the effects of
the monitor, solver rotation, and elite transfer. A factorial ablation with and
without memory, adaptive triggering, and portfolio diversity is needed for
component-level attribution.

\subsection{Statistical reliability and remaining validity threats}
\label{sec:discussion-validity}

The inferential protocol strengthens the conclusions in several ways. Each
configuration used 31 stochastic runs; non-normal distributions were handled
with rank-based tests; paired designs were used where common seeds were
documented; and Holm correction controlled family-wise error within the stated
comparison families. Conclusions were synthesized using normalized gaps,
within-problem ranks, and W/T/L decisions rather than by pooling incompatible
raw objective scales. This is essential because fitness in MKP, error in
CEC2022, and LCOE in HRES2--H2 are not directly commensurable.

Several limitations remain. First, the test coverage is finite: MKP uses the
first instance from nine Chu--Beasley families, CEC2022 is evaluated only at
$D=10$, and HRES2--H2 represents one synthetic annual scenario. Second, the
present parameter study compares alignment mode and MKP budget, but it does not
systematically vary the monitoring-window length, Sakoe--Chiba band,
percentiles, no-improvement threshold, or patience. Third, raw per-run baseline
arrays were not archived for MKP and HRES2--H2. The reported HRES2--H2
Mann--Whitney probabilities and Holm arithmetic can be checked, but the base
statistics and effect sizes cannot be independently reconstructed from the
available summaries. The HRES2--H2 hybrid--baseline samples also lack the
common-seed pairing available for the direct DTW--DDTW comparison. Fourth, the
analysis reports significance and directional summaries but does not include
confidence intervals for paired differences or standardized non-parametric
effect sizes. Finally, post hoc selection of the best observed configuration
should not be treated as validation on unseen instances.

These limitations motivate a clear validation agenda: archive all run-level
outputs; preregister comparison families and practically meaningful margins;
report effect sizes and bootstrap confidence intervals; test additional MKP
instances and CEC dimensions; evaluate HRES2--H2 under multiple annual and
economic scenarios; and conduct controlled component ablations. An external
test set or nested configuration protocol would further distinguish genuine
generalization from benchmark-specific tuning.

\subsection{Practical implications}
\label{sec:discussion-implications}

For the tested MKP setting, DTW--3K is the most defensible default when the
additional computational budget is acceptable. For CEC2022 and HRES2--H2, the
current evidence does not justify a universal preference between DTW and DDTW;
the simpler level-based DTW variant is competitive, while DDTW remains useful
when trajectory shape is expected to be more informative than absolute level.
In either case, selection should be validated on representative instances and
should account for the cost of objective evaluation.

Overall, the strongest conclusion is that trajectory-triggered portfolio
rotation can provide reliable cross-domain performance without requiring a
single standalone metaheuristic to be optimal everywhere. The evidence is most
decisive for the budgeted MKP comparison and for the hybrid's distributional
advantage in the specific HRES2--H2 case; it is more nuanced on CEC2022, where
the framework is the best aggregate performer but specialized methods remain
competitive on individual functions. This balance between robustness and
specialization should define the contribution: the framework is an adaptive
portfolio mechanism with broad empirical strength, not a universal dominance
guarantee.
